% Numerical Method section - Algorithm-focused version
\section{Numerical Method}
\label{sec:method}

The numerical method combines the Allen-Cahn phase-field model for interface tracking with the lattice Boltzmann method for hydrodynamics, following the pressure-velocity formulation of Fakhari \& Bolster \cite{fakhari2017diffuse}. This section first describes the governing equations and their LBM discretization, then focuses on the tensor-based implementation details that distinguish this work from traditional CUDA implementations.

\subsection{Allen-Cahn Phase-Field Model}

We employ the conservative Allen-Cahn equation to track the liquid-gas interface. The order parameter $\phi \in [0,1]$ transitions smoothly from $\phi = 1$ (liquid) to $\phi = 0$ (gas) over a characteristic thickness $\xi$. The free energy of the system comprises a double-well bulk potential and a gradient penalty:

\begin{equation}
    F = \int_\Omega \left[ \beta \phi^2(1-\phi)^2 + \frac{\kappa}{2} |\nabla \phi|^2 \right] d\Omega,
    \label{eq:free_energy}
\end{equation}

where $\beta$ is the bulk free energy parameter and $\kappa$ is the gradient energy coefficient. The chemical potential is:

\begin{equation}
    \mu = \frac{\delta F}{\delta \phi} = 4\beta \phi(\phi-1)(\phi-0.5) - \kappa \nabla^2 \phi.
    \label{eq:chemical_potential}
\end{equation}

The surface tension is related to these parameters by:
\begin{equation}
    \sigma = \frac{\sqrt{2\kappa\beta}}{6}.
    \label{eq:surface_tension}
\end{equation}

The conservative Allen-Cahn equation governing interface evolution is:
\begin{equation}
    \frac{\partial \phi}{\partial t} + \nabla \cdot (\phi \mathbf{u}) = M_\phi \nabla^2 \mu,
    \label{eq:ac_conservative}
\end{equation}
where $M_\phi$ is the mobility. In the sharp-interface limit ($\xi \to 0$), this reduces to the standard Allen-Cahn equation with sharpening.

The mobility is related to the diffusion coefficient through $M_\phi = M/\beta$, where $M$ is typically set to $0.02/\beta$ to maintain a balance between interface diffusion and sharpening.

\subsection{Lattice Boltzmann Formulation}

We use a two-distribution approach with D3Q19 lattice structure. The $h$-distribution tracks the phase field:
\begin{equation}
    h_\alpha (\mathbf{x} + \mathbf{e}_\alpha \Delta t, t + \Delta t) = h_\alpha (\mathbf{x}, t) - \frac{1}{\tau_h} \left( h_\alpha - h_\alpha^{eq} \right) + S_\alpha^h,
    \label{eq:h_evolution}
\end{equation}
with equilibrium:
\begin{equation}
    h_\alpha^{eq} = w_\alpha \phi \left( 1 + \frac{\mathbf{e}_\alpha \cdot \mathbf{u}}{c_s^2} \right).
    \label{eq:h_equilibrium}
\end{equation}

The source term $S_\alpha^h$ incorporates the Allen-Cahn sharpening force \cite{fakhari2017diffuse}.

The $g$-distribution solves the incompressible Navier-Stokes equations:
\begin{equation}
    g_\alpha (\mathbf{x} + \mathbf{e}_\alpha \Delta t, t + \Delta t) = g_\alpha (\mathbf{x}, t) - \frac{1}{\tau_g} \left( g_\alpha - g_\alpha^{eq} \right) + S_\alpha^g,
    \label{eq:g_evolution}
\end{equation}
with equilibrium distribution of Fakhari \& Bolster:
\begin{equation}
    g_\alpha^{eq} = w_\alpha \left[ P + \rho c_s^2 \left( \frac{\mathbf{e}_\alpha \cdot \mathbf{u}}{c_s^2} + \frac{(\mathbf{e}_\alpha \cdot \mathbf{u})^2}{2c_s^4} - \frac{\mathbf{u} \cdot \mathbf{u}}{2c_s^2} \right) \right].
    \label{eq:g_equilibrium}
\end{equation}

The source term $S_\alpha^g$ includes surface tension, pressure correction, and viscous correction forces.

\textbf{Density-dependent relaxation time} ensures stability at high density ratios:
\begin{equation}
    \tau_{\text{eff}} = \frac{1}{2} + \frac{\phi \rho_l (\tau_l - \tfrac{1}{2}) + (1-\phi) \rho_g (\tau_g - \tfrac{1}{2})}{\phi \rho_l + (1-\phi) \rho_g}.
    \label{eq:tau_eff}
\end{equation}

Macroscopic quantities are recovered as:
\begin{align}
    \phi &= \sum_\alpha h_\alpha, \quad
    \rho \mathbf{u} = \frac{1}{c_s^2} \sum_\alpha \mathbf{e}_\alpha g_\alpha - \frac{\Delta t}{2} \phi \nabla \mu.
    \label{eq:macro_recover}
\end{align}

\subsection{Volume Penalization for Curved Boundaries}

Traditional bounce-back boundary conditions on Cartesian grids produce staircase approximations to curved geometries. Volume penalization (VP) addresses this by introducing a continuous solid fraction field $\chi(\mathbf{x}) \in [0,1]$ \cite{song2022volume}:

\begin{equation}
    \chi(\mathbf{x}) = 0.5 \left[ 1 + \tanh\left( \frac{2(R_{\text{ridge}} - r)}{\xi} \right) \right],
    \label{eq:chi_smoothing}
\end{equation}
where $r$ is the distance to the ridge center. The transition occurs over approximately one interface thickness $\xi$, providing smooth boundary normals for contact angle calculations.

During the collision step, VP modulates the distribution functions:
\begin{equation}
    h_\alpha^{\text{post}} = (1 - \chi) h_\alpha^{\text{stream}} + \chi \cdot h_{\bar{\alpha}}^{\text{pre}},
    \label{eq:vp_collision}
\end{equation}
eliminating staircase artifacts while maintaining mass conservation.

\subsection{Geometric Gradient Amplification}

Achieving extreme contact angles ($\theta_{eq} = 162^\circ$) requires overcoming the resistance of the AC sharpening term. The geometric gradient amplification technique enhances the wall-normal interface gradient:

\begin{equation}
    (\nabla \phi)_{\text{corrected}} = \nabla \phi + \text{amp} \times \frac{\mu_{\text{wall}}}{|\nabla \phi|} \mathbf{n}_\phi,
    \label{eq:amp_correction}
\end{equation}
where $\mu_{\text{wall}}$ is determined by Young's equation:
\begin{equation}
    \cos \theta_{eq} = -\frac{\phi_c}{\sqrt{2\kappa\beta}}.
    \label{eq:young_equation}
\end{equation}

The amplification coefficient follows approximately $\text{amp} \propto N^{0.84}$, where $N$ is the grid resolution.

% ==============================================================================
% NEW SECTION: Tensor-Based Implementation Details
% ==============================================================================
\subsection{Tensor-Based Implementation in PyTorch}

This subsection describes the algorithmic choices that distinguish this implementation from traditional CUDA codes.

\subsubsection{Tensor Layout and Memory Organization}

The LBM state is organized as contiguous tensors on the GPU. For a D3Q19 lattice with grid shape $(N_x, N_y, N_z)$, the distribution functions are stored with the direction dimension last, enabling efficient broadcasting during collision and streaming:

\begin{verbatim}
# Tensor layout: (nx, ny, nz, Q) for D3Q19
# Example: h.shape = (150, 150, 107, 19) for 150^3 grid
h = torch.zeros((nx, ny, nz, Q), dtype=torch.float32, device='cuda')
\end{verbatim}

This layout ensures that the innermost dimension (Q) maps to contiguous GPU memory, maximizing memory coalescing during the collision step where all Q directions are accessed per grid node.

\textbf{Critical memory optimization}: Instead of cloning full distribution fields for bounce-back, we clone only the solid-node distributions. For a typical simulation with 20\% solid fraction on a $150^3$ grid, this reduces memory usage by approximately 15\%.

\subsubsection{Streaming via \texttt{torch.roll}}

The LBM streaming step $\mathbf{f}(\mathbf{x}, t+\Delta t) = \mathbf{f}(\mathbf{x} - \mathbf{e}_\alpha \Delta t, t)$ is implemented using \texttt{torch.roll}:

\begin{verbatim}
for i in range(1, Q):
    h[..., i] = torch.roll(h_pre[..., i], shifts=-e[i], dims=d)
\end{verbatim}

This approach leverages PyTorch's fused CUDA kernel for roll operations, avoiding explicit Python loops over lattice directions. The broadcast-based indexing also enables vectorized gradient computations across all directions simultaneously.

\subsubsection{Isotropic Gradient Computation}

Gradients of the phase field are computed using the lattice weights as coefficients, ensuring isotropy on the D3Q19 lattice:

\begin{verbatim}
grad = torch.zeros((ndim, nx, ny, nz), dtype=torch.float32, device='cuda')
for i in range(1, Q):
    for d in range(ndim):
        grad[d] += w[i] * e[i][d] * shifted_phi / cs2
\end{verbatim}

This tensor-based gradient operator avoids the separate kernel launches required in traditional implementations and enables fusion with subsequent operations.

\subsubsection{Memory Management and Tensor Lifecycle}

Intermediate tensors are deleted immediately after use to minimize peak memory usage:

\begin{verbatim}
# Before: del intermediate_tensor after use
mu_phi = 4 * beta * phi * (phi - 1) * (phi - 0.5) - kappa * lap_phi
del lap_phi  # Free immediately
\end{verbatim}

The floating-point format is fixed at \texttt{float32} for production runs, as \texttt{float64} approximately doubles memory usage and reduces throughput by 30-40\% on NVIDIA GPUs without tensor cores.

\subsubsection{Multi-Backend Support}

The implementation supports multiple GPU backends through PyTorch's device abstraction:

\begin{itemize}
    \item \textbf{NVIDIA CUDA}: Standard \texttt{device='cuda'}, supports mixed precision via \texttt{torch.cuda.amp}.
    \item \textbf{Intel XPU}: Via \texttt{device='xpu'} with \texttt{pip install torch --index-url https://download.pytorch.org/whl/xpu}.
    \item \textbf{CPU}: Fallback for debugging (\texttt{device='cpu'}).
\end{itemize}

The solver handles non-standard devices by falling back to raw device objects when \texttt{torch.device} instantiation fails, enabling seamless execution across hardware without code changes.

\begin{table}[htbp]
    \centering
    \caption{Simulation parameters (lattice units) for standard configuration.}
    \label{tab:params}
    \begin{tabular}{lcc}
        \toprule
        Parameter & Symbol & Value \\
        \midrule
        Liquid density & $\rho_l$ & 828.0 \\
        Gas density & $\rho_g$ & 1.0 \\
        Density ratio & $\rho_l/\rho_g$ & 828:1 \\
        Droplet diameter & $D_0$ & 45 (150$^3$) \\
        Interface thickness & $\xi$ & 5.0 \\
        Liquid relaxation time & $\tau_l$ & 0.53 \\
        Gas relaxation time & $\tau_g$ & 0.60 \\
        Equilibrium contact angle & $\theta_{eq}$ & $162^\circ$ \\
        Impact velocity & $U_0$ & $-0.05$ \\
        Grid resolution & $N^3$ & $100^3$ / $150^3$ \\
        \bottomrule
    \end{tabular}
\end{table}
